% =====================================================================
% 第 V 章  Class fields (K. Iwasawa, Nov. 20 and 27, 1957)
% =====================================================================
\lecture{Class fields}{Class Fields}{K. Iwasawa}{Nov. 20 and 27, 1957}

% ---------------------------------------------------------------
\section{Introduction}

As was proved in IV, §5, the class invariants $j(\mathfrak{k})$ of an
imaginary quadratic field $K$ satisfy certain fundamental congruences. In
the following, we shall first give an outline of class field theory in
its classical form and then show, using that theory, how we can deduce
from those congruences arithmetic properties of the numbers
$j(\mathfrak{k})$ and of the extension $K(j(\mathfrak{k}))/K$. For the
details of classical class field theory, cf.\ Hasse's
\emph{Klassenk\"orper Bericht} quoted in I, §4 \cite{hasse1926bericht}.

% ---------------------------------------------------------------
\section{Ideal groups}

Let $K$ be a finite algebraic number field. Let $m$ be a divisor of $K$,
i.e.\ a formal product of a finite number of prime divisors of $K$:
$m=\prod\mathfrak{p}_{i}^{e_{i}}$. Put $m=m_{0}m_{\infty}$, where $m_{0}$
and $m_{\infty}$ are, respectively, products of non-archimedean and
archimedean prime divisors in $m=\prod\mathfrak{p}_{i}^{e_{i}}$; $m_{0}$
may be then identified with the corresponding ideal of $K$. In the
following, we shall always consider such a divisor $m$ that $m_{0}$ is an
integral ideal of $K$ and that $m_{\infty}$ is a product of a number of
distinct real archimedean prime divisors of $K$. For a number $\xi$ in
$K$, we then write
\[
\xi\equiv 1\bmod m
\]
if and only if $\nu_{\mathfrak{p}_{i}}(\xi-1)\geq e_{i}$ for every
$\mathfrak{p}_{i}$ dividing $m_{0}$ and $\sigma_{\mathfrak{p}_{j}}(\xi)>0$
for every $\mathfrak{p}_{j}$ dividing $m_{\infty}$; here
$\nu_{\mathfrak{p}_{i}}$ denotes the normalized exponential valuation of
$K$ belonging to $\mathfrak{p}_{i}$ and $\sigma_{\mathfrak{p}_{j}}$
denotes the isomorphism of $K$ into the real field corresponding to
$\mathfrak{p}_{j}$.

Now, let $I=I(K)$ be the multiplicative group of all non-zero ideals of
$K$. For any divisor $m$ as considered above, we denote by $I_{m}
=I_{m}(K)$ the subgroup of $I$ consisting of all ideals of $K$ which are
prime to $m_{0}$ and by $S_{m}=S_{m}(K)$ the group of all principal
ideals $(\xi)$ with $\xi\equiv 1\bmod m$. $S_{m}$ is then a subgroup of
$I_{m}$ and $I_{m}/S_{m}$ is a finite group. Any group $H$ such that
$S_{m}\subset H\subset I_{m}$ will be called an \emph{ideal group} of $K$
defined mod $m$.

% ---------------------------------------------------------------
\section{The density of a set of prime ideals}

Let $P=P(K)$ be the set of all prime ideals of $K$. For any subset $M$ of
$P$, we put
\[
\mu(s;M)=\sum_{\mathfrak{p}\in M}N_{K/\mathbb{Q}}(\mathfrak{p})^{-s}.
\]
The right hand side is absolutely convergent for $\operatorname{Re}(s)>1$
and $\mu(s;M)$ is well defined in that domain. Now, if the limit
\[
\delta(M)=\lim_{s>1,\ s\to 1}\frac{\mu(s;M)}{\mu(s;P)}
\]
exists, we call $\delta(M)$ the \emph{density} (Dirichlet density or
Kronecker density) of the set $M$ in $P$. Clearly, $\delta(P)=1$.

Using the fact that the zeta-function
$\zeta_{K}(s)=\prod_{\mathfrak{p}\in P}
\bigl(1-N_{K/\mathbb{Q}}(\mathfrak{p})^{-s}\bigr)^{-1}$ has a simple pole
at $s=1$, we can see by a simple computation that
\[
\mu(s;P)\sim\log\zeta_{K}(s)\sim\log\frac{1}{s-1},
\]
where $\sim$ indicates that the difference of the both sides of $\sim$ is
a function bounded at $s=1$. Hence
\[
\lim_{s\to 1}\mu(s;P)\big/\log\frac{1}{s-1}=1,
\]
and the density $\delta(M)$ can be defined also by
\[
\delta(M)=\lim_{s\to 1}\mu(s;M)\big/\log\frac{1}{s-1}.
\]

\emph{Example.} Let $N_{K/\mathbb{Q}}(\mathfrak{p})=p^{d}$ ($p$ rational
prime), $d=\deg\mathfrak{p}$ = absolute degree of $\mathfrak{p}$. Let
$P'=P'(K)$ be the set of all $\mathfrak{p}$ with $d=1$ and let
$P''=P''(K)$ be the complement of $P'$ in $P$. Since
$N_{K/\mathbb{Q}}(\mathfrak{p})=p^{d}\geq p^{2}$ for $\mathfrak{p}$ in
$P''$, $\mu(s;P'')$ is bounded at $s=1$. Hence,
$\lim\mu(s;P'')/\log\frac{1}{s-1}=0$, and
\[
\delta(P'')=0,\qquad\delta(P')=\delta(P)=1.
\]

More generally, if we put $M'=M\cap P'$, $M''=M\cap P''$ for any subset
$M$ of $P$, then
\[
\delta(M'')=0,\qquad\delta(M')=\delta(M),
\]
whenever one of $\delta(M')$ and $\delta(M)$ exists.

Now, let $H$ be an ideal group of $K$ defined mod $m$ and let
$\mathfrak{k}$ be any class (coset) of $I_{m}/H$. For the set
$\mathfrak{k}\cap P$, we have then
\[
\delta(\mathfrak{k}\cap P)=\delta(\mathfrak{k}\cap P')
=\frac{1}{h},
\tag{1}\label{eq:V-1}
\]
where $h=[I_{m}:H]$. This is the theorem of arithmetic progression for
the field $K$, which generalizes the well-known theorem of Dirichlet for
$K=\mathbb{Q}$.

An outline of the proof of E.~Hecke for \eqref{eq:V-1} is as follows
(cf.\ G\"ottinger Nachrichten, 1917): Let
\[
L(s;\chi)=\prod'\bigl(1-\chi(\mathfrak{p})
N_{K/\mathbb{Q}}(\mathfrak{p})^{-s}\bigr)^{-1},
\qquad\operatorname{Re}(s)>1,
\]
be the L-function of $K$ defined for a character $\chi$ of $I_{m}/H$
($\prod'$ over $\mathfrak{p}$ in $I_{m}$).

It is first proved that $L(s;\chi)$ can be analytically continued to a
meromorphic function of $s$ on the entire s-plane which has a unique pole
of order $1$ at $s=1$ if $\chi=\chi_{0}$ and is holomorphic everywhere if
$\chi\neq\chi_{0}$; here $\chi_{0}$ denotes the principal character of
$I_{m}/H$. On the other hand, it follows easily from the definition that
\[
\prod_{\chi}L(s;\chi)
=\prod'\bigl(1-N_{K/\mathbb{Q}}(\mathfrak{p})^{-fs}\bigr)^{-h/f},
\qquad\operatorname{Re}(s)>1,
\]
where $\prod_{\chi}$ is taken over all characters of $I_{m}/H$ and
$f=f(\mathfrak{p})$ is the order of the class of $\mathfrak{p}$ in
$I_{m}/H$. Hence,
\[
\prod_{\chi}L(s;\chi)=\sum_{\nu=1}^{\infty}\frac{a_{\nu}}{\nu^{s}},
\qquad\operatorname{Re}(s)>1,
\]
with rational integers $a_{\nu}\geq 0$. Suppose now that $L(1;\chi)=0$
for some $\chi\neq\chi_{0}$. Then the left hand side of the above would
be holomorphic on the entire s-plane and, by a theorem of Landau on the
convergence of Dirichlet series with positive coefficients, the right
hand side of the above must converge for all $s$. Since the $a_{\nu}$ are
integers, it would then follow that $a_{\nu}=0$ except for a finite
number -- a contradiction. Hence:
\[
L(1;\chi)\neq 0,\qquad\chi\neq\chi_{0}.
\]

Now, by a simple computation, we have
\[
\begin{aligned}
\log L(s;\chi)
&=\Sigma'\sum_{\nu=1}^{\infty}\frac{1}{\nu}\,
\chi(\mathfrak{p})^{\nu}N_{K/\mathbb{Q}}(\mathfrak{p})^{-\nu s}\\
&\sim\Sigma'\,\chi(\mathfrak{p})\,N_{K/\mathbb{Q}}(\mathfrak{p})^{-s}
\end{aligned}
\qquad(s>1),
\]
By the orthogonality of the characters $\chi$, it then follows that
\[
\sum_{\chi}\chi(\mathfrak{k}^{-1})\log L(s;\chi)
\sim h\cdot\sum_{\mathfrak{p}\in\mathfrak{k}}
N_{K/\mathbb{Q}}(\mathfrak{p})^{-s}
=h\cdot\mu(s;\mathfrak{k}\cap P),
\]
for any class $\mathfrak{k}$ of $I_{m}/H$. Since $L(1;\chi)\neq 0$ for
$\chi\neq\chi_{0}$, we have
\[
\log L(s;\chi_{0})\sim\log\frac{1}{s-1};\qquad
\log L(s;\chi)\sim 0,\quad\chi\neq\chi_{0},
\]
and the left hand side of the above is $\sim\log\frac{1}{s-1}$. Hence we
get
\[
\delta(\mathfrak{k}\cap P)=\lim_{s\to 1}
\mu(s;\mathfrak{k}\cap P)\big/\log\frac{1}{s-1}=\frac{1}{h},
\]
q.e.d.

\subsection*{Remark.}
By Theorems II, IV of §6 below (which can be proved purely
algebraically), there exists an abelian extension $E$ of degree $h$ over
$K$ such that the zeta-function $\zeta_{E}(s)$ of $E$ is the product of
$\prod_{\chi}L(s;\chi)$ and a finite number of factors of the form
$(1-q^{-s})^{-1}$, $q>1$. Since $\zeta_{E}(s)$ has a simple pole at
$s=1$, this immediately implies $L(1;\chi)\neq 0$ for $\chi\neq\chi_{0}$.
For $K=\mathbb{Q}$, such an $E$ is given by a subfield of a cyclotomic
field.

% ---------------------------------------------------------------
\section{The inequality h <= n}\label{sec:V-inequality}

Let $L$ be a Galois extension of degree $n$ over $K$. A prime ideal
$\mathfrak{p}$ of $K$ is said to be completely decomposed in $L$ if
$\mathfrak{p}$ is decomposed into the product of $n$ distinct prime
ideals of $L$, or equivalently, if $\mathfrak{p}$ is unramified in $L$
and $\mathfrak{p}=N_{L/K}(\mathfrak{P})$ for some prime ideal
$\mathfrak{P}$ of $L$. Let $W$ denote the set of all such
$\mathfrak{p}$ of $K$ which are completely decomposed in $L$, and let
$P''(L)$ be the set of all prime ideals $\mathfrak{P}$ of $L$ such that
$\deg\mathfrak{P}=1$ and that $\mathfrak{p}=N_{L/K}(\mathfrak{P})$ is
unramified in $L$. Then
$\mathfrak{P}\longrightarrow\mathfrak{p}=N_{L/K}(\mathfrak{P})$ defines
an $n$-to-$1$ correspondence between the sets $P''(L)$ and
$W'=W\cap P'(K)$ and, as
$N_{L/\mathbb{Q}}(\mathfrak{P})=N_{K/\mathbb{Q}}(\mathfrak{p})$, we get
$\mu(s;P''(L))=n\,\mu(s;W')$.

However, the difference of $P'(L)$ and $P''(L)$ is a finite set and
$\delta(P''(L))=\delta(P'(L))=\delta(P(L))=1$. Hence, it follows from the
above that
\[
\delta(W)=\delta(W')=\frac{1}{n}.
\tag{2}\label{eq:V-2}
\]

Now, let $m$ be a divisor of $K$ as considered in §2. Define a group
$N_{m}=N_{m}(L/K)$ by
\[
N_{m}=S_{m}(K)\,N_{L/K}\bigl(I_{m}(L)\bigr).
\]
Clearly, $S_{m}\subset N_{m}\subset I_{m}$ and we put
\[
h_{m}=h_{m}(L/K)=[I_{m}:N_{m}].
\]
By the definition of $W$, $W_{m}=W\cap I_{m}$ is contained in
$N_{L/K}(I_{m}(L))$ and, hence, also in $N_{m}$. As the difference of $W$
and $W_{m}$ is again a finite set, we obtain from \eqref{eq:V-1} and
\eqref{eq:V-2} that
\[
\frac{1}{n}=\delta(W)=\delta(W_{m})
\leq\delta(N_{m}\cap P)=\frac{1}{h_{m}},
\]
namely that
\[
h_{m}\leq n,\qquad\text{or}\qquad
[I_{m}:N_{m}(L/K)]\leq[L:K].
\tag{3}\label{eq:V-3}
\]
This is called the second fundamental inequality of class field theory
and it is valid for any $m$ and for any finite Galois extension $L$ of
$K$.

% ---------------------------------------------------------------
\section{Definition of class fields}

Now, a finite Galois extension $L$ of a finite algebraic number field $K$
is called a \emph{class field} over $K$ if, in the above, $h_{m}=n$ holds
for some divisor $m$ of $K$. If $L/K$ is such a class field, there exists
a unique divisor $f$ of $K$ such that $h_{m}=n$ holds if and only if $m$
is a multiple of $f$. The divisor $f=f(L/K)$ is then called the
\emph{conductor} of $L/K$.

\begin{namedstmt}{Criterion}
If $L/K$ is a class field, there exists a divisor $m$ of $K$ and an ideal
group $H$ of $K$ defined mod $m$ such that

i) $N_{L/K}(I_{m}(L))\subset H$;

ii) $H\cap P$ is contained in $W$ except for a subset of density $0$.

Conversely, if there exist such $m$ and $H$, then $L/K$ is a class field
and
\[
h_{m}=n,\qquad H=N_{m}.
\]
\end{namedstmt}

\begin{proof}[Proof.]
If $L/K$ is a class field, put $H=N_{m}$ for a divisor $m$ such that
$h_{m}=n$. Then i) holds obviously and ii) follows from that the density
of the difference of $N_{m}\cap P$ and $W_{m}$ is $0$. Suppose,
conversely, such $m$ and $H$ exist. By i), we have
$S_{m}\subset N_{m}\subset H\subset I_{m}$. Put $a=[H:N_{m}]$. Then, by
the theorem of arithmetic progression, $\delta(H\cap P)=\frac{a}{h_{m}}$.
However, by ii), $\delta(H\cap P)\leq\delta(W)=\frac{1}{n}$. Hence
$\frac{a}{h_{m}}\leq\frac{1}{n}$, or, $an\leq h_{m}$. It then follows
from \eqref{eq:V-3} that $a=1$, $h_{m}=n$, q.e.d.
\end{proof}

As an immediate application we notice the following: Let both
$L_{1}/K$ and $L_{2}/K$ be class fields and let $m$ be divisible by
$f(L_{1}/K)$ and $f(L_{2}/K)$. It is then easy to see that $m$ and
$H=N_{m}(L_{1}/K)\cap N_{m}(L_{2}/K)$ satisfy the above conditions i) and
ii) for the extension $L_{1}L_{2}/K$. Hence $L_{1}L_{2}/K$ is also a
class field and $h_{m}(L_{1}L_{2}/K)=[L_{1}L_{2}:K]$,
$N_{m}(L_{1}L_{2}/K)=N_{m}(L_{1}/K)\cap N_{m}(L_{2}/K)$. Now, the
condition $N_{m}(L_{1}/K)\subset N_{m}(L_{2}/K)$ is obviously equivalent
with the condition $N_{m}(L_{1}L_{2}/K)=N_{m}(L_{1}/K)$ and, hence, also
with $[L_{1}L_{2}:K]=h_{m}(L_{1}L_{2}/K)=h_{m}(L_{1}/K)=[L_{1}:K]$, or,
$L_{1}L_{2}=L_{1}$. Therefore, $L_{1}\supset L_{2}$ if and only if
$N_{m}(L_{1}/K)\subset N_{m}(L_{2}/K)$.

% ---------------------------------------------------------------
\section{Fundamental theorems}\label{sec:V-fundamental}

We are now going to state the fundamental theorems of class field theory.

\begin{namedstmt}{I}
A finite extension $L$ of $K$ is a class field over $K$ if and only if
$L/K$ is an abelian extension (i.e.\ a Galois extension with an abelian
Galois group).
\end{namedstmt}

The most essential step of the proof is to prove
\[
n\leq h_{m}
\]
for some $m$, when $L/K$ is a cyclic extension. This is called the first
fundamental inequality of class field theory.

\begin{namedstmt}{II}
Given any ideal group $H$ of $K$ defined mod $m$, there exists a unique
class field $L$ over $K$ such that $[I_{m}:N_{m}(L/K)]=[L:K]$ and
$H=N_{m}(L/K)$.
\end{namedstmt}

\begin{namedstmt}{III}
A prime divisor $\mathfrak{p}$ of $K$ is ramified in a class field $L$
over $K$ if and only if $\mathfrak{p}$ divides the conductor $f(L/K)$ of
$L/K$.
\end{namedstmt}

Before going to state the next theorem, Artin's reciprocity law, we first
make some preparations. Let $L/K$ be an arbitrary Galois extension of
degree $n$ with the Galois group $G$. Let $\mathfrak{p}$ be a prime ideal
of $K$ unramified in $L$, and $\mathfrak{P}$ a prime ideal of $L$
dividing $\mathfrak{p}$. Then there exists a unique element $\sigma$ in
$G$ such that
\[
\sigma(\alpha)\equiv\alpha^{N_{K/\mathbb{Q}}(\mathfrak{p})}
\bmod\mathfrak{P}
\tag{4}\label{eq:V-4}
\]
for any algebraic integer $\alpha$ in $L$. This element $\sigma$ is called
the \emph{Frobenius substitution} of $\mathfrak{P}$ for the extension
$L/K$ and is denoted by $\sigma_{\mathfrak{P}}$; it generates the
so-called decomposition group of the prime ideal $\mathfrak{P}$, and if
$f$ is the order of $\sigma_{\mathfrak{P}}$ and $n=fg$, then
\[
N_{L/K}(\mathfrak{P})=\mathfrak{p}^{f},\qquad
\mathfrak{p}=\mathfrak{P}_{1},\dots,\mathfrak{P}_{g},\qquad
\mathfrak{P}_{1}=\mathfrak{P},
\]
where the $\mathfrak{P}_{i}$ are distinct prime ideals of $L$.

Now, assume that $L/K$ is a class field, i.e.\ an abelian extension. Then
the Frobenius substitutions of $\mathfrak{P}_{1},\cdots,\mathfrak{P}_{g}$
coincide with each other and we may denote these
$\sigma_{\mathfrak{P}_{i}}$ $(i=1,\ldots,g)$ simply by
$\sigma_{\mathfrak{p}}$. Let $m$ be a divisor of $K$ such that $h_{m}=n$
and let $\mathfrak{a}=\prod_{i=1}^{r}\mathfrak{p}_{i}^{e_{i}}$ be the
prime ideal decomposition of an ideal $\mathfrak{a}$ in $I_{m}$. By III,
every $\mathfrak{p}_{i}$ is unramified in $L$ and $\sigma_{\mathfrak{p}_{i}}$
is defined as above. So, we may define $\sigma_{\mathfrak{a}}$ by
\[
\sigma_{\mathfrak{a}}=\prod_{i=1}^{r}\sigma_{\mathfrak{p}_{i}}^{e_{i}}.
\]
Clearly, $\mathfrak{a}\longrightarrow\sigma_{\mathfrak{a}}$ is a
homomorphism of $I_{m}$ into $G$. But we have the following theorem:

\begin{namedstmt}{IV}
The homomorphism $\mathfrak{a}\longrightarrow\sigma_{\mathfrak{a}}$
induces an isomorphism
\[
I_{m}/N_{m}\cong G;
\]
namely, the homomorphism is surjective and the kernel is $N_{m}$.
\end{namedstmt}

An immediate consequence of IV is the following: Prime ideals of $K$,
which are contained in the same class of $I_{m}/N_{m}$, are decomposed
into the product of the same number of prime factors in $L$. In
particular, every prime ideal in $N_{m}$ is completely decomposed in $L$
without exception.

Finally, let $L/K$ be again a Galois extension of degree $n$ and let
$\sigma$ be an arbitrary element of the Galois group $G$ of $L/K$. We
denote by $P_{\sigma}$ the set of all prime ideals $\mathfrak{p}$ of $K$
such that $\mathfrak{p}$ is unramified in $L$ and that
$\sigma=\sigma_{\mathfrak{p}}$ for some prime ideal $\mathfrak{P}$ of $L$
dividing $\mathfrak{p}$. Then

\begin{namedstmt}{V}
If $c$ is the number of elements in the conjugate class of $\sigma$ in
$G$, then
\[
\delta(P_{\sigma})=\frac{c}{n}.
\]
\end{namedstmt}

Notice that if $L/K$ is a class field, V is an immediate consequence of
IV and the theorem of arithmetic progression.

% ---------------------------------------------------------------
\section{Hilbert's absolute class field}

Let $1$ denote the unit divisor of $K$. Then $I_{1}(K)=I(K)$, and
$S_{1}(K)$ is the group of all principal ideals of $K$. The factor group
$I_{1}/S_{1}$ is hence the ideal class group $C_{K}$ of $K$ in the
ordinary sense and its order $h$ is nothing but the class number of $K$.

Now, by the existence theorem II, there exists a class field $L_{0}$ over
$K$ such that
\[
h=[I_{1}:S_{1}]=[L_{0}:K],\qquad S_{1}=N_{1}(L_{0}/K).
\]
The conductor $f(L_{0}/K)$ is then $1$ and, by III, $L_{0}/K$ is an
unramified extension. Take an arbitrary unramified abelian extension
$L/K$. Again by III, the conductor $f(L/K)$ is $1$, and since
$S_{1}=N_{1}(L_{0}/K)$ is contained in $N_{1}(L/K)$, $L$ is contained in
$L_{0}$ (cf.\ the last lines of §5). Thus we see that $L_{0}$ is the
unique maximal unramified abelian extension of $K$. $L_{0}$ is called the
\emph{Hilbert's absolute class field} over $K$. By IV, the Galois group
of $L_{0}/K$ is canonically isomorphic with the ideal class group
$C_{K}=I_{1}/S_{1}$ of $K$.

% ---------------------------------------------------------------
\section{Conjugacy of class invariants}\label{sec:V-conjugacy}

We now assume that $K$ is an imaginary quadratic field and denote by
$\mathfrak{k}_{1},\ldots,\mathfrak{k}_{n}$ the ideal classes of
$C_{K}=I_{1}/S_{1}$. The class invariants $j(\mathfrak{k}_{i})$,
$i=1,\ldots,h$, are then distinct algebraic integers. Take a finite Galois
extension $L$ over $K$ containing all such class invariants
$j(\mathfrak{k}_{i})$ and denote by $\xi_{1},\ldots,\xi_{s}$ the distinct
elements in the set $\sigma(j(\mathfrak{k}_{i}))$ where $\sigma$ runs over
the Galois group $G$ of $L/K$ and $1\leq i\leq h$. We then put
\[
\gamma=\prod_{a<b}(\xi_{a}-\xi_{b})
\]
and denote by $P^{*}$ the set of all prime ideals $\mathfrak{p}$ of $K$
with $\deg\mathfrak{p}=1$ such that $\mathfrak{p}$ is unramified in $L$
and is prime to $\gamma$. It is clear that
$\delta(P^{*})=\delta(P')=1$. Hence, for any subset $M$ of $P(K)$ with
$\delta(M)>0$, $\delta(M\cap P^{*})=\delta(M)>0$ and the intersection
$M\cap P^{*}$ is non-empty.

Let $\mathfrak{p}$ be a prime ideal in $P^{*}$, $\mathfrak{k}_{p}$ the
ideal class of $C_{K}$ containing $\mathfrak{p}$, and
$p=N_{K/\mathbb{Q}}(\mathfrak{p})$. By IV, §5, we have then
\[
j(\mathfrak{k}_{p}^{-1}\mathfrak{k})\equiv j(\mathfrak{k})^{p}
\bmod\mathfrak{p},
\]
for any ideal class $\mathfrak{k}$. On the other hand, if $\mathfrak{P}$
is a prime ideal of $L$ dividing $\mathfrak{p}$ and $\sigma_{\mathfrak{P}}$
is the Frobenius substitution of $\mathfrak{P}$ for $L/K$, then by
\eqref{eq:V-4},
\[
\sigma_{\mathfrak{P}}(j(\mathfrak{k}))\equiv j(\mathfrak{k})^{p}
\bmod\mathfrak{P}.
\]
It follows that
\[
\sigma_{\mathfrak{P}}(j(\mathfrak{k}))
\equiv j(\mathfrak{k}_{\mathfrak{p}}^{-1}\mathfrak{k})
\bmod\mathfrak{P}.
\]
But, the both sides of the above congruence are numbers in the set
$\xi_{1},\ldots,\xi_{s}$. Hence, as $\mathfrak{p}$ is prime to $\gamma$,
we get
\[
\sigma_{\mathfrak{P}}(j(\mathfrak{k}))
=j(\mathfrak{k}_{p}^{-1}\mathfrak{k}),
\tag{5}\label{eq:V-5}
\]
\[
\mathfrak{p}\in P^{*},\qquad\mathfrak{P}\mid\mathfrak{p}.
\]

Now, let $\mathfrak{k}$ and $\mathfrak{k}'$ be arbitrary ideal classes of
$C_{K}$. By the theorem of arithmetic progression,
$\delta(\mathfrak{k}\mathfrak{k}'^{-1}\cap P^{*})
=\delta(\mathfrak{k}\mathfrak{k}'^{-1}\cap P)>0$ and the ideal class
$\mathfrak{k}\mathfrak{k}'^{-1}$ contains a prime ideal $\mathfrak{p}$ in
$P^{*}$ so that $\mathfrak{k}'=\mathfrak{k}_{p}^{-1}\mathfrak{k}$. It then
follows from \eqref{eq:V-5} that the class invariants $j(\mathfrak{k})$
and $j(\mathfrak{k}')$ are conjugate in $L$ over $K$. On the other hand,
given any $\sigma$ in $G$, $\delta(P_{\sigma}\cap P^{*})
=\delta(P_{\sigma})>0$ by V and there exists a prime ideal
$\mathfrak{p}$ in $P^{*}$ such that $\sigma=\sigma_{\mathfrak{P}}$ for
some $\mathfrak{P}$ dividing $\mathfrak{p}$. Hence, by \eqref{eq:V-5},
any conjugate $\sigma(j(\mathfrak{k}))$ of a class invariant
$j(\mathfrak{k})$ is another class invariant $j(\mathfrak{k}')$. Thus:

\setcounter{theorem}{0}
\begin{theorem}\label{thm:V-1}
The class invariants $j(\mathfrak{k}_{i})$, $i=1,\ldots,h$, of an
imaginary quadratic field $K$ form a complete set of conjugates over $K$.
\end{theorem}

% ---------------------------------------------------------------
\section{The extension K(j(k)) / K}\label{sec:V-extension}

From now on, we put $L=K(j(\mathfrak{k}_{1}),\ldots,j(\mathfrak{k}_{n}))$,
for the latter is a Galois extension of $K$ by
Theorem~\ref{thm:V-1}.

\begin{namedstmt}{Lemma}
A prime ideal $\mathfrak{p}$ in $P^{*}$ is completely decomposed in $L$
if and only if $\mathfrak{p}$ is contained in $S_{1}$ (i.e.\ a principal
ideal).
\end{namedstmt}

\begin{proof}[Proof.]
Let $\mathfrak{P}$ be a prime ideal of $L$ dividing $\mathfrak{p}$. Then
$\mathfrak{p}$ is completely decomposed in $L$ if and only if
$\sigma_{\mathfrak{P}}=1$. Since
$L=K(j(\mathfrak{k}_{1}),\ldots,j(\mathfrak{k}_{n}))$, it follows from
\eqref{eq:V-5} that $\sigma_{\mathfrak{P}}=1$ if and only if
$j(\mathfrak{k}_{p}^{-1}\mathfrak{k})=j(\mathfrak{k})$ for every
$\mathfrak{k}$, namely, if and only if $\mathfrak{k}_{p}=1$, or,
$\mathfrak{p}\in S_{1}$.
\end{proof}

We now verify that the two conditions of the criterion given in §5 are
satisfied for $m=1$, $H=S_{1}$ and for the extension $L/K$. Let
$\mathcal{O}$ be an arbitrary ideal of $L$. By the theorem of arithmetic
progression, there exists a prime ideal $\mathfrak{P}$ of $L$ such that
$\deg\mathfrak{P}=1$ and $\mathcal{O}=(\alpha)\mathfrak{P}$ with a number
$\alpha$ in $L$. We may also assume that
$\mathfrak{p}=N_{L/K}(\mathfrak{P})$ is prime to $\gamma$. Then
$\mathfrak{p}$ is contained in $P^{*}$ and is completely decomposed in
$L$. By the above lemma, $\mathfrak{p}$ is hence contained in $S_{1}$,
and so is the ideal
$N_{L/K}(\mathcal{O})=N_{L/K}(\alpha)N_{L/K}(\mathfrak{P})$. Thus, i) is
verified. As $\delta(S_{1}\cap P^{*})=\delta(S_{1}\cap P)$, the condition
ii) is also satisfied by the above lemma.

By the criterion, $L/K$ is therefore a class field and
\[
h=[L:K],\qquad S_{1}=N_{1}(L/K).
\]
As explained in §7, $L$ is hence the Hilbert's absolute class field over
$K$, i.e.\ the maximal unramified abelian extension of $K$. Since $L/K$
is now abelian and $j(\mathfrak{k}_{i})$, $1\leq i\leq h$, are conjugate
over $K$, we can obtain $L$ by adjoining just one $j(\mathfrak{k}_{i})$ to
$K$. Thus the following theorem is proved:

\begin{theorem}\label{thm:V-2}
Let $j(\mathfrak{k})$ be any class invariant of an imaginary quadratic
field $K$. Then $K(j(\mathfrak{k}))$ is the maximal unramified abelian
extension of $K$.
\end{theorem}

As mentioned in §7, we have then a canonical isomorphism between the
ideal class group $C_{K}$ of $K$ and the Galois group $G$ of $L/K$:
\[
C_{K}\cong G.
\]
Let $\sigma_{\mathfrak{k}}$ denote the element of $G$ corresponding to an
ideal class $\mathfrak{k}$ under this isomorphism. If $\mathfrak{p}$ is a
prime ideal of $P^{*}$ contained in the class $\mathfrak{k}$, then
$\sigma_{\mathfrak{k}}=\sigma_{\mathfrak{p}}$ by the definition of the
isomorphism. We get therefore from \eqref{eq:V-5} the following explicit
formula for the automorphism $\sigma_{\mathfrak{k}}$ of $L/K$:
\[
\sigma_{\mathfrak{k}}(j(\mathfrak{k}'))
=j(\mathfrak{k}^{-1}\mathfrak{k}'),
\tag{6}\label{eq:V-6}
\]
\[
\mathfrak{k},\mathfrak{k}'\in C_{K}.
\]

Now, let $\tau$ be an automorphism of the field of all algebraic numbers.
Clearly, $\tau(L)$ is the maximal unramified abelian extension of
$\tau(K)$. Since $K$ is a quadratic field over $\mathbb{Q}$ and
$\tau(K)=K$, it follows that $\tau(L)=L$. Therefore, $L/\mathbb{Q}$ is
also a Galois extension. Let $\mathcal{G}$ denote the Galois group of
$L/\mathbb{Q}$ and let $\rho$ be the element of $\mathcal{G}$ which maps
$\alpha$ in $L$ to $\bar{\alpha}$ (conjugate complex of $\alpha$).
Clearly, $\rho^{2}=1$. But $\rho$ induces a non-trivial automorphism on
the imaginary field $K$. Hence the order of $\rho$ is $2$ and
$\mathcal{G}$ is the semi-direct product of the subgroup $\{1,\rho\}$ and
the normal subgroup $G$. Let $\mathfrak{a}$ be an ideal in an ideal class
$\mathfrak{k}$. Then
$\rho(j(\mathfrak{k}))=\overline{j(\mathfrak{k})}
=j(\overline{\mathfrak{k}})$. Since
$\mathfrak{a}\bar{\mathfrak{a}}=N_{K/\mathbb{Q}}(\mathfrak{a})$ is a
principal ideal, $\overline{\mathfrak{k}}$ is in the class
$\mathfrak{k}^{-1}$ and we have
\[
\rho(j(\mathfrak{k}))=j(\mathfrak{k}^{-1}).
\tag{7}\label{eq:V-7}
\]

Combined with Theorem~\ref{thm:V-1}, this gives us the following:

\begin{namedstmt}{THEOREM 1$'$}
The class invariants $j(\mathfrak{k}_{i})$, $i=1,\ldots,h$, of an
imaginary quadratic field $K$ form a complete set of conjugates over
$\mathbb{Q}$.
\end{namedstmt}

Therefore, if we put $f(t)=\prod_{i}\bigl(t-j(\mathfrak{k}_{i})\bigr)$,
$f(t)$ is an irreducible polynomial of $t$ with integral rational
coefficients.

Another consequence of \eqref{eq:V-7} is the following: Let $\sigma$ be
an arbitrary element of $G$ and let $\sigma=\sigma_{\mathfrak{k}}$,
$\sigma^{-1}=\sigma_{\mathfrak{k}^{-1}}$. Since $G$ is normal in
$\mathcal{G}$, $\rho\sigma\rho^{-1}$ is again contained in $G$, and it
follows from \eqref{eq:V-6} and \eqref{eq:V-7} that
\[
(\rho\sigma\rho^{-1})(j(\mathfrak{k}'))
=\rho\sigma\bigl(j(\mathfrak{k}'^{-1})\bigr)
=\rho\bigl(j(\mathfrak{k}^{-1}\mathfrak{k}'^{-1})\bigr)
=j(\mathfrak{k}\mathfrak{k}')
=\sigma^{-1}(j(\mathfrak{k}')),
\]
for any class invariant $j(\mathfrak{k}')$. Hence
\[
\rho\sigma\rho^{-1}=\sigma^{-1},\qquad\sigma\in G.
\]
This, together with the fact that $\mathcal{G}$ is the semi-direct product
of $\{1,\rho\}$ and $G$, shows that the structure of the Galois group
$\mathcal{G}$ is completely determined when we know the structure of
$G\cong C_{K}$.
